% GENERATED FILE. Edit canonical lessons, then run npm run generate:books.
% canonical-source-hash: fnv1a64:73eaeca7738bf485
% canonical-lessons: PT-C460-javali, PT-C460-veado, PT-C460-lebre, PT-C460-falcao, PT-C460-pastar

\chapter{Animais selvagens}
\label{ch:pt-animais-selvagens}
\hlchaptermodality{460}

\begin{chapteropening}
\textbf{By the end of this chapter:} \emph{I can name wild animals I might see: the wild boar, the deer, the hare, the falcon, and animals grazing.}

Complete the last lesson of chapter 460: I can name wild animals I might see: the wild boar, the deer, the hare, the falcon, and animals grazing.
\end{chapteropening}

\section[o javali]{o javali — the wild boar}

\label{lesson:PT-C460-javali}



\begin{quote}
Before the new one: say the Portuguese for to tan, then the Portuguese for to dive.
\end{quote}

\subsection*{You'll want to know: o javali}
\textbf{o javali} — \textquotedblleft{}the wild boar\textquotedblright{}.

\begin{quote}
À noite, os javalis descem até às hortas.
\end{quote}

\begin{quote}
\emph{At night, the wild boar come down to the vegetable gardens.}
\end{quote}

Feminine: \emph{a javalina}.

\begin{cousinweb}
\emph{javali} comes from Andalusian Arabic \emph{jabalī} 'wild boar' (EtymDB), literally 'of the mountain', short for 'mountain pig'.
\end{cousinweb}

\begin{grammarlens}[title={What you've built}]
One more word for telling, explaining and describing at B1.
\end{grammarlens}

\subsection*{Your turn}
\begin{itemize}
\raggedright
  \item \emph{Say it:} \emph{o javali}
  \item \emph{Say it:} \emph{o javali}, once more
  \item \emph{Say it:} a sentence of your own with \emph{o javali}
  \item \emph{Recall:} say the Portuguese for to tan, then the Portuguese for to dive, then say \emph{o javali} again
\end{itemize}

\begin{tcolorbox}[breakable,colback=teal!4,colframe=teal!35!black,title={Before you move on}]
What does o javali mean? (\textquotedblleft{}The wild boar\textquotedblright{}.) Say it once more.
\end{tcolorbox}
\section[o veado]{o veado — the deer}

\label{lesson:PT-C460-veado}



\begin{quote}
Before the new one: say the Portuguese for to dive, then the Portuguese for the wild boar.
\end{quote}

\subsection*{You'll want to know: o veado}
\textbf{o veado} — \textquotedblleft{}the deer\textquotedblright{}.

\begin{quote}
Vimos um veado à beira da estrada.
\end{quote}

\begin{quote}
\emph{We saw a deer at the side of the road.}
\end{quote}

\emph{O cervo} is the red deer; \emph{veado} covers most deer.

\begin{cousinweb}
\emph{veado} is inherited through Old Portuguese \emph{veado} (EtymDB) from Latin \emph{vēnātus} 'hunting, game' (Lewis \& Short): the animal people hunted.
\end{cousinweb}

\begin{grammarlens}[title={What you've built}]
One more word for telling, explaining and describing at B1.
\end{grammarlens}

\subsection*{Your turn}
\begin{itemize}
\raggedright
  \item \emph{Say it:} \emph{o veado}
  \item \emph{Say it:} \emph{o veado}, once more
  \item \emph{Say it:} a sentence of your own with \emph{o veado}
  \item \emph{Recall:} say the Portuguese for to dive, then the Portuguese for the wild boar, then say \emph{o veado} again
\end{itemize}

\begin{tcolorbox}[breakable,colback=teal!4,colframe=teal!35!black,title={Before you move on}]
What does o veado mean? (\textquotedblleft{}The deer\textquotedblright{}.) Say it once more.
\end{tcolorbox}
\section[a lebre]{a lebre — the hare}

\label{lesson:PT-C460-lebre}



\begin{quote}
Before the new one: say the Portuguese for the wild boar, then the Portuguese for the deer.
\end{quote}

\subsection*{You'll want to know: a lebre}
\textbf{a lebre} — \textquotedblleft{}the hare\textquotedblright{}.

\begin{quote}
A lebre corre mais depressa do que o coelho.
\end{quote}

\begin{quote}
\emph{The hare runs faster than the rabbit.}
\end{quote}

\emph{Gato por lebre} — a cat for a hare: being cheated.

\begin{cousinweb}
\emph{lebre} is inherited from Latin \emph{leporem}, the accusative of \emph{lepus} 'hare' (EtymDB).
\end{cousinweb}

\begin{grammarlens}[title={What you've built}]
One more word for telling, explaining and describing at B1.
\end{grammarlens}

\subsection*{Your turn}
\begin{itemize}
\raggedright
  \item \emph{Say it:} \emph{a lebre}
  \item \emph{Say it:} \emph{a lebre}, once more
  \item \emph{Say it:} a sentence of your own with \emph{a lebre}
  \item \emph{Recall:} say the Portuguese for the wild boar, then the Portuguese for the deer, then say \emph{a lebre} again
\end{itemize}

\begin{tcolorbox}[breakable,colback=teal!4,colframe=teal!35!black,title={Before you move on}]
What does a lebre mean? (\textquotedblleft{}The hare\textquotedblright{}.) Say it once more.
\end{tcolorbox}
\section[o falcão]{o falcão — the falcon}

\label{lesson:PT-C460-falcao}



\begin{quote}
Before the new one: say the Portuguese for the deer, then the Portuguese for the hare.
\end{quote}

\subsection*{You'll want to know: o falcão}
\textbf{o falcão} — \textquotedblleft{}the falcon\textquotedblright{}.

\begin{quote}
Um falcão voa por cima do campo.
\end{quote}

\begin{quote}
\emph{A falcon is flying over the field.}
\end{quote}

Very fast when it dives.

\begin{cousinweb}
\emph{falcão} is inherited from Late Latin \emph{falcō, falcōnis} 'falcon' (EtymDB), often linked with \emph{falx} 'sickle' for its curved claws.
\end{cousinweb}

\begin{grammarlens}[title={What you've built}]
One more word for telling, explaining and describing at B1.
\end{grammarlens}

\subsection*{Your turn}
\begin{itemize}
\raggedright
  \item \emph{Say it:} \emph{o falcão}
  \item \emph{Say it:} \emph{o falcão}, once more
  \item \emph{Say it:} a sentence of your own with \emph{o falcão}
  \item \emph{Recall:} say the Portuguese for the deer, then the Portuguese for the hare, then say \emph{o falcão} again
\end{itemize}

\begin{tcolorbox}[breakable,colback=teal!4,colframe=teal!35!black,title={Before you move on}]
What does o falcão mean? (\textquotedblleft{}The falcon\textquotedblright{}.) Say it once more.
\end{tcolorbox}
\section[pastar]{pastar — to graze}

\label{lesson:PT-C460-pastar}



\begin{quote}
Before the new one: say the Portuguese for the hare, then the Portuguese for the falcon.
\end{quote}

\subsection*{You'll want to know: pastar}
\textbf{pastar} — \textquotedblleft{}to graze\textquotedblright{}.

\begin{quote}
As vacas pastam no prado.
\end{quote}

\begin{quote}
\emph{The cows graze in the meadow.}
\end{quote}

\emph{O pasto} — the pasture.

\begin{cousinweb}
\emph{pastar} is built on \emph{pasto} 'pasture' (EtymDB), which is inherited from Latin \emph{pāstus} 'feeding, pasture', from \emph{pāscere} 'to feed'.
\end{cousinweb}

\begin{grammarlens}[title={What you've built}]
One more word for telling, explaining and describing at B1.
\end{grammarlens}

\subsection*{Your turn}
\begin{itemize}
\raggedright
  \item \emph{Say it:} \emph{pastar}
  \item \emph{Say it:} \emph{pastar}, once more
  \item \emph{Say it:} a sentence of your own with \emph{pastar}
  \item \emph{Recall:} say the Portuguese for the hare, then the Portuguese for the falcon, then say \emph{pastar} again
\end{itemize}

\begin{tcolorbox}[breakable,colback=teal!4,colframe=teal!35!black,title={Before you move on}]
What does pastar mean? (\textquotedblleft{}To graze\textquotedblright{}.) Say it once more.
\end{tcolorbox}
